Desastres naturais provocados por eventos climáticos extremos geram impactos profundos. No âmbito econômico, recursos
destinados à recuperação de áreas afetadas enfraquecem a microeconomia local e ampliam a vulnerabilidade futura. No
entanto, o impacto social costuma ser mais severo. Traumas, perdas humanas e rupturas comunitárias deixam danos
duradouros, muitas vezes irreversíveis.

Devido às mudanças climáticas, a frequência e intensidade dos desastres naturais vêm aumentando globalmente, como
resultado, muitas regiões são afetadas por eventos climáticos extremos de maneira acelerada, causando um significativo
despreparo. Entre esses desastres, a ocorrência de enchentes se destacam pelo aumento de extremos de precipitação
globais, como evidenciado por Dunn et al. (2024) \cite{dunn}. As mesmas apresentam uma capacidade de causar destruição
generalizada e afetar áreas densamente povoadas, especialmente em regiões urbanas com infraestrutura inadequada para
gerenciar grandes volumes de água. Em muitos casos, essa inadequação é agravada pela falta de manutenção de sistemas
essenciais, como drenagem urbana ou canais fluviais, tornando as cidades vulneráveis mesmo quando dispõem de estruturas
projetadas para lidar com a situação.

No final de abril e início de maio de 2024, o Rio Grande do Sul registrou precipitações excepcionais que elevaram
drasticamente os níveis das bacias hidrográficas. Na região metropolitana de Porto Alegre (RMPOA), ocorreu a maior cheia
da história do Rio Guaíba, superando o recorde da Cheia de 1941 \cite{dep}. Em 5 de maio, o nível atingiu 5{,}33 metros,
e falhas no sistema de contenção - muros, diques, comportas e estações de bombeamento \cite{comercio} -, resultaram no
alagamento de bairros inteiros, evacuações nas áreas mais afetadas, bloqueio das principais vias e interrupções de
abastecimento, energia e mobilidade. O evento decorreu de fatores meteorológicos combinados à vulnerabilidade urbana
\cite{yang, hammond, zhang, pereira}, deixando milhares de deslocados e centenas de vítimas \cite{alcantara}.

Durante a crise, previsões confiáveis do nível do Guaíba foram essenciais para orientar ações emergenciais. Os modelos
hidrológicos e hidrodinâmicos utilizados, embora precisos, exigem grande volume de dados topográficos, meteorológicos e
hidrológicos, conhecimento especializado e tempo de simulação elevado. Nos momentos críticos, só puderam ser executados
diariamente \cite{iph}, limitando sua utilidade em previsões de curto prazo \cite{atashi}.

Modelos baseados em dados surgem como alternativa mais ágil. Eles utilizam técnicas para identificar e explorar relações
estatísticas entre os dados de entrada e saída. Com a aprendizagem dos padrões presentes nos dados históricos, é
possível mapear as dinâmicas subjacentes do sistema e utilizá-las para prever valores futuros. Assim, esses modelos são
capazes de realizar previsões precisas com base em séries temporais, capturando tendências e comportamentos complexos
que podem ser difíceis de modelar diretamente por meio de abordagens tradicionais. Enquanto esses modelos demandam muito
tempo e dados específicos, os modelos baseados em dados oferecem maior agilidade para previsões em tempo real, sendo uma
escolha mais adequada para sistemas de alerta precoce e mitigação de desastres de cheias \cite{ta}.

Entre os métodos clássicos, destaca-se o ARIMA \cite{box}, que combina autoregressão, média móvel e diferenciação, além
de permitir sazonalidade via SARIMA e variáveis exôgenas através do SARIMAx. Esses modelos têm histórico de aplicação em
hidrologia, como Atashi et al.\ (2022) \cite{atashi} e Güldal et al.\ (2009) \cite{gudal}, servindo como base de
comparação para métodos modernos.

No campo de \textit{Machine Learning}, técnicas de \textit{Gradient Boosting} são amplamente utilizadas pela rapidez e
precisão, como o XGBoost \cite{chen_xgboost} e o LightGBM \cite{lightgbm}, que mantêm desempenho elevado com menor custo
computacional. Em \textit{Deep Learning}, modelos recorrentes como LSTM \cite{hochreiter} capturam dependências
temporais complexas e têm se mostrado eficazes em predições hidrológicas \cite{borwarnginn}.

Este estudo aplica essas abordagens à previsão do nível da água durante as enchentes do Guaíba, avaliando a eficácia de
SARIMA, XGBoost, LSTM e outros métodos em um cenário crítico, onde precisão e rapidez são essenciais. A comparação será
feita por métricas consolidadas, como \textit{Kling–Gupta efficiency} (KGE), \textit{Nash-Sutcliffe Model Efficiency
Coefficient} (NSE), \textit{Root Mean Squared Error} (RMSE) e \textit{Mean Absolute Error} (MAE), permitindo avaliar
tanto a acurácia quanto a capacidade de cada modelo de reproduzir a dinâmica hidrológica.

Também será analisado o uso de variáveis externas, como precipitação, vento e afluentes, além de aplicar transformações
como \textit{lag}, janelas temporais e cumulativos. O objetivo é identificar não apenas o menor erro preditivo, mas o
comportamento de cada técnica frente às diferentes dinâmicas do sistema. Essa análise contribui para estratégias mais
robustas de gestão de desastres.

Para esta análise, serão utilizados dados  obtidos das estações telemétricas relevantes através do Sistema Nacional de
Informações sobre Recursos Hídricos (SNIRH), que disponibiliza informações das principais bacias hidrográficas do Brasil
\cite{snirh}. Esses dados serão a base para as séries históricas do nível das águas do Guaíba, complementados por dados
externos, como precipitação, vazão e temperatura também do SNIRH.